\documentclass{report}
\usepackage{amsmath,amssymb}
\setlength{\parindent}{0mm}
\setlength{\parskip}{1em}
\begin{document}
\begin{center}
\rule{6in}{1pt} \
{\large Daniel Szyld \\
{\bf Asynchronous Optimized Schwarz Methods for the solution of PDEs}}

Temple University \\ Department of Mathematics \\ 1805 N Broad St \\ Philadelphia PA 19122
\\
{\tt szyld@temple.edu}\\
Fr\'ederic Magoul\'es\\
Cedric Venet\end{center}

Asynchronous methods are parallel iterative procedures where each process
performs its task without waiting for other processes to be completed.
For the numerical solution of a general PDE on a domain,
Schwarz iterative methods use a decomposition
into two or more (overlapping) subdomains.
In the classical formulation, Dirichlet boundary conditions are used
on the artificial interfaces.
Given an initial approximation, the method
progresses by solving the PDE in
each subdomain using as boundary data on the artificial interfaces
the values of the solution on the neighboring subdomain from the previous step.
This procedure is inherently parallel.
For optimized Schwarz, the boundary conditions on the artificial interfaces are
of Robin or mixed type.
Thus, one can optimize the Robin parameter(s) obtaining a very fast method.

We present an asynchronous version of the optimized Schwarz method
for the solution of differential equations on a parallel computational environment.
In a one-way subdivision of the computational domain, with overlap,
the method is shown to converge
when the optimal artificial interface conditions are used.
Convergence is also proved under very mild conditions on the size of the subdomains,
when approximate (non-optimal) interface conditions are utilized.
Numerical results are presented on large three-dimensional problems
illustrating the efficiency of the proposed
asynchronous parallel implementation of the method.


\end{document}
